% REVISED 2026-06-07 compression: ~3pp -> ~1.5pp, descriptive-texture recitation moved to App H
% REVISED: canonical-target reframe 2026-06-04
% REVISED: hostile-review action plan 2026-06-04
%   M3a: §5.1 re-sequenced to LEAD with the opportunity-adjustment
%   vanish (raw SMDs demoted to descriptive texture); §5.2 subsection
%   folded into §5.1; referee 5.2 explicit non-CADE-anchor sentence
%   added; referee 5.4 "partly mechanical = not independent
%   corroboration" deflation moved to the lead; closing summary
%   reconciled (no claimed network specificity). No numbers changed.
%
% Award-Layer Screening" reframe. Strengths-forward, honest about
% attenuation, standardized differences (not p-values from large
% samples). Object is adjudication-anchored exposure, not membership.
%
% five main tables and both figures deleted; full profile detail now
% lives in Appendix H (\ref{app:profile_submission}); refs repointed.
\section{Economic Content and Ordinary-Loser Alternatives}
\label{sec:cobidder_type}

Validation establishes that the award-layer ranking concentrates
adjudication-anchored loser-side exposure on the $\valMainCobidders$
always-loser cobidders, of which $\valMainCobFL$ ($52.4\%$) satisfy the
frequent-loser cut and $\valMainCobNonFL$ ($47.6\%$) do not. What that
target lacks so far is economic content. We supply it through a
demanding within-stratum comparison: holding the cut fixed (so no gap
can be read back into it), do the $\valMainCobFL$ FL cobidders look
different from FL non-cobidders along dimensions that would matter for
loser-side prioritization? With thousands of firms per group $p$-values
carry no information, so we report standardized mean differences (SMD,
Cohen's $d$), which speak to magnitude. The full profile, the
opportunity-adjusted differences, the monotonicity bins, the
binary-versus-continuous comparison, and the ordinary-loser
adjudication table appear in \ref{app:profile_submission}; here we keep
to the headline findings.

\subsection{What opportunity adjustment removes, what survives, and what the rank orders}
\label{sec:profile_5_1}

The headline cuts against the screen. Once we adjust for
opportunity---the expected contact and exposure a firm would
accumulate given its participation profile---the breadth, persistence,
and specialization gaps between cobidders and other frequent losers
\emph{vanish}. The item-group concentration difference falls from a raw
SMD of $\valProfHHIigRawSMD$ to an adjusted $\valProfHHIigAdjSMD$, so
the apparent product specialization of cobidders is a feature of the
procurement environments they compete in, not an intrinsic portfolio
trait; the breadth (distinct buyers) and persistence (active years)
gaps go the same way. The raw contrasts that vanish---a higher Preg\~ao
share, more distinct buyers, modestly longer active spans---are
descriptive texture rather than findings, recited in full in
\ref{app:profile_submission}. What survives is only the mechanical
defendant-contact variables: proximity to the legal anchors attenuates
modestly, with the CADE-cases SMD slipping from $\valProfNcasesSMD$ to
$\valProfNcasesAdjSMD$. That residual is no corroboration, because the
label is itself built from defendant contact: the mechanical share
cannot be peeled apart from a behavioral one, so the contact variables
cannot stand as evidence about conduct.

Network specificity points the same way. Negative controls show the
raw concentration is generic co-participation arithmetic rather than a
signature of the adjudicated network: swapping the true CADE defendants
for matched placebo anchors reproduces the raw ranking (real-anchor
ROC-AUC $\valProfNegCtrlRealAUC$ against a placebo-anchor mean of
$\valProfNegCtrlPlaceboAUC$; $p=0.46$), and anchors built from non-CADE
high-volume winners---who sit at least as tightly inside the
network---discriminate at least as well (null mean $0.78$; $p=0.91$).
The global zero-win definition does retain more raw discrimination than
any market-specific alternative (best alternative
$\valProfMarketZeroWinBest$ against the global $\valProfNegCtrlRealAUC$),
so the definition itself is not where the construct weakens
(\ref{app:profile_submission}).

The most telling diagnostic concerns what the ranking actually orders.
Binning firms by participation $T_i$, cobidder prevalence climbs
monotonically with participation---from near zero to roughly $28\%$ in
the highest bin (rank correlation $\valProfMonoPrevRho$)---while
opportunity-adjusted excess contact $X_i = O_i - E_i$ stays
flat-to-negative across the same bins (rank correlation
$\valProfMonoExcessRho$). Higher-ranked firms are not contacting
defendants \emph{more than their exposure predicts}; they are simply
more exposed. The award-layer rank is therefore an \textbf{exposure
ranking, not a collusion-intensity ranking}. This confirms the FL14
cutoff as administrative rather than structural: the continuous score
$\log(1+T_i)$ ties or beats the binary in every full-sample
configuration (full-universe ROC-AUC $0.761$ continuous vs.\ $0.688$
binary), and there is no bunching at the threshold---the density ratio
at $T=14$ is $\valProfBunchRatio$, essentially flat
(\ref{app:profile_submission}). The reading is consistent with evidence
that auction collusion can operate through bidder roles, repeated
relationships, and cartel organization
\citep{baldwin1997bidder,porter1999ohio,pesendorfer2000study,asker2010leniency}.

\subsection{Ordinary-loser alternatives and what this section establishes}
\label{sec:ordinary_5_4}

Could the same profile arise from the ordinary, non-collusive reasons
firms lose persistently? Ordinary-loser alternatives are \emph{narrowed
but not eliminated}. Cobidders are wider and longer-lived rather than
narrow and short-lived, and carry a lower Convite and minimum-quorum
share, weakening the low-capacity quorum-filler story without ruling it
out; but the proxies that would settle the rest---distance-to-winner,
geography and logistics, later wins---are unobserved at the award layer,
so benign specialization and geography stay on the table
(\ref{app:profile_submission}). The one axis these stories do not
predict is the one the screen is built on: contact with the adjudicated
CADE network.

On balance the limits dominate. Most of the raw distinction is
procurement-environment composition and partly mechanical proximity
that vanishes once opportunity is held fixed; the raw ranking is
reproduced by placebo and non-CADE anchors; the rank is monotone in
exposure yet flat in opportunity-adjusted excess contact; FL14 is
administrative rather than structural; and ordinary-loser alternatives
are narrowed, not eliminated. The construct is a disciplined exposure
ranking whose non-mechanical economic content is modest. It establishes
no membership, agreement, liability, or cover bidding in any tender,
and on its own it is not diagnostic. Read alongside the validation of
Section~\ref{sec:validation} and the sequencing of
Section~\ref{sec:screening_forensics}, it supports triage rather than
liability: the award layer orders attention toward
adjudication-anchored loser-side exposure, and richer records then
judge what the screen has prioritized.
Section~\ref{sec:screening_forensics} turns to implementation, asking
whether the award-layer rank can be placed ahead of bid-distribution
forensics in a sequence that spares bid-level reconstruction while
preserving adjudication-anchored recall.
